\section{四槽异步 FIFO：多 bit 数据怎样跨越两个时钟域}

两级同步器适合单 bit 状态，却不能保证多 bit 数据来自同一个源时刻。异步 FIFO 把数据与跨域控制分开：数据写入双口存储，写侧只推进写指针，读侧只推进读指针；跨时钟域传递的是采用格雷码编码的指针副本。接收方看到旧指针只会暂时保守地判断“仍空”或“仍满”，不会读取尚未发布的数据，也不会提前覆盖尚未归还的槽位。

下面使用深度为 4 的 FIFO。存储槽位为 0--3，读写指针各使用 3 bit：低 2 bit 选择槽位，最高 bit 记录是否绕环。二进制指针便于本时钟域加一和寻址；格雷指针由
\[
  G=B\mathbin{\oplus}(B\gg1)
\]
得到，相邻值只改变 1 bit，适合逐位经过同步器。二进制指针不直接跨域。

\begin{pdfdiagram}[x=1cm,y=1cm]
  \node[hw control, minimum width=2.8cm] (wctrl) at (1.6,4.3) {写时钟域\\写使能与二进制写指针};
  \node[hw memory, minimum width=3.0cm] (mem) at (6.3,4.3) {四槽双口存储\\写口与读口独立};
  \node[hw control, minimum width=2.8cm] (rctrl) at (11.0,4.3) {读时钟域\\读使能与二进制读指针};
  \draw[hw bus] (wctrl) -- node[above, hw label]{写数据/地址} (mem);
  \draw[hw bus] (mem) -- node[above, hw label]{读数据} (rctrl);

  \node[hw state, minimum width=2.8cm] (wgray) at (1.6,2.1) {写指针格雷码\\只由写域产生};
  \node[hw state, minimum width=2.8cm] (wseen) at (11.0,2.1) {同步后的写指针\\读域判断 empty};
  \draw[hw bus] (wgray) -- node[above, hw label]{两级同步，允许晚看到} (wseen);

  \node[hw state, minimum width=2.8cm] (rseen) at (1.6,0.2) {同步后的读指针\\写域判断 full};
  \node[hw state, minimum width=2.8cm] (rgray) at (11.0,0.2) {读指针格雷码\\只由读域产生};
  \draw[hw return] (rgray) -- node[above, hw label]{两级同步，允许晚看到} (rseen);
\end{pdfdiagram}
\diagramnote{异步 FIFO 把数据面和控制面分开。上排数据只经过双口存储；中、下两排分别传递格雷写指针和格雷读指针，方向不同但路线分离。每个时钟域只修改自己的指针，并用同步到本域的远端指针生成空满状态。}

\topic{四次写入把 FIFO 填满}

初始读写指针均为二进制 \code{000}、格雷码 \code{000}，FIFO 为空。写侧连续写入 \code{3A}、\code{7C}、\code{55}、\code{E1}。每次有效写在同一个写时钟沿完成两件事：把数据写入当前低 2 bit 指定的槽位，并把写指针推进到下一位置。第四次写后，写指针从 \code{011} 绕到 \code{100}，格雷码从 \code{010} 变为 \code{110}；它与尚未推进的读指针相隔整整 4 个槽位，因此写侧置 full，不再接受第五次写。

\begin{longtable}{L{0.10\linewidth}L{0.17\linewidth}L{0.16\linewidth}L{0.18\linewidth}L{0.27\linewidth}}
\toprule
事件 & 写入槽位/数据 & 写指针 $B\to G$ & 写侧状态 & 读侧何时可见 \\
\midrule
复位完成 & 无 & \code{000 -> 000} & empty，非 full & 两域先确认共同空闲代际 \\
$W_0$ & 槽 0/\code{3A} & \code{001 -> 001} & 占用 1 槽 & 格雷指针同步后才解除 empty \\\tablerowrule
$W_1$ & 槽 1/\code{7C} & \code{010 -> 011} & 占用 2 槽 & 可能仍只看到 $W_0$ \\\tablerowrule
$W_2$ & 槽 2/\code{55} & \code{011 -> 010} & 占用 3 槽 & 旧指针只造成保守延迟 \\\tablerowrule
$W_3$ & 槽 3/\code{E1} & \code{100 -> 110} & full，拒绝新写 & 最终看到四个已发布槽位 \\
\bottomrule
\end{longtable}

深度为 4 时，指针多出的最高 bit 区分“同一槽位编号但相差一整圈”。empty 的判断简单：读侧的下一格雷读指针等于同步后的格雷写指针。full 的判断需要识别相差一圈：写侧的下一格雷写指针等于同步读指针在格雷表示下翻转最高两位后的结果。具体比较电路由深度决定，但语义始终是“下一次写是否会追上仍未归还的最旧数据”。

\topic{读取释放槽位，但写侧不会立即知道}

读侧看到非 empty 后，在自己的时钟沿从槽 0 取出 \code{3A}，再把读指针从 \code{000} 推进到 \code{001}。槽 0 的所有权已经在读域归还，但新读指针还要编码为格雷码并经过两级同步，写域才会解除 full。同步延迟期间，写侧宁可多拒绝一两个周期，也不能猜测槽 0 已经安全可写。

这种保守性是异步 FIFO 正确性的核心。远端指针可能旧，但不能凭空超前：读侧不会在数据发布之前看到写指针前进，写侧也不会在数据消费之前看到读指针前进。因此旧状态只降低瞬时吞吐，不会破坏数据所有权。

\begin{longtable}{L{0.14\linewidth}L{0.26\linewidth}L{0.30\linewidth}L{0.20\linewidth}}
\toprule
边界 & 本域已经发生的事实 & 远端尚未看到的事实 & 安全后果 \\
\midrule
写提交 & 存储槽写完，写指针同沿推进 & 读域同步器仍保存旧写指针 & 读侧暂时保持 empty，不会读早 \\\tablerowrule
读提交 & 数据已交给读侧，读指针同沿推进 & 写域同步器仍保存旧读指针 & 写侧暂时保持 full，不会覆盖早 \\\tablerowrule
同步完成 & 新格雷指针到达远端第二级 & 远端可重新计算空或满 & 只改变接收/拒绝决策，不搬数据 \\\tablerowrule
非法请求 & full 时写或 empty 时读 & 指针与存储均保持 & 数据顺序和槽位所有权不变 \\
\bottomrule
\end{longtable}

\topic{独立复位需要协议化重新对齐}

如果两个时钟域同时进入复位，清零两侧指针即可得到共同空 FIFO。困难发生在单侧复位：若读域把指针清零，而写域仍保留旧指针和旧数据，双方可能把同一组格雷码解释成新的空满关系；若写域单独清零，读域还可能把旧槽位当作新代际数据。仅清寄存器不足以证明跨域所有权已经一致。

支持独立复位的系统需要在 FIFO 外增加恢复协议：先停止新写和新读，等待或丢弃在途事务；两侧各自进入 reset-done 状态；同步交换代际号或重新初始化标志；只有双方都观察到同一代际且指针为零，才重新开放 valid/ready。旧代际数据可以被明确丢弃，但不能在没有身份区分的情况下混入新流。

故障定位也应沿边界进行。数据内容错误而指针顺序正确，先查双口存储写使能、地址和数据宽度；重复或跳项，查指针推进条件是否与真正握手同沿；偶发 full/empty 抖动，查格雷码转换、同步级和组合跨域；复位后才出错，查两域释放顺序与代际握手。异步 FIFO 的完成不是“数据写进数组”，而是写指针发布、读侧同步观察、读提交以及槽位归还全部按各自时钟边界闭合。
